\documentclass[conference]{IEEEtran}
\IEEEoverridecommandlockouts

% ============================================================
% PACKAGES
% ============================================================
\usepackage{cite}
\usepackage{amsmath,amssymb,amsfonts}
\usepackage{graphicx}
\usepackage{textcomp}
\usepackage{xcolor}
\usepackage{booktabs}
\usepackage{array}
\usepackage{bm}
\usepackage{url}
\usepackage{multirow}
\graphicspath{{figures/}} % Put manuscript figures in the figures/ directory.

% ============================================================
% MASTER-DRAFT NOTE -- FROZEN HRI RESEARCH CONTRACT v1.1
% Research identity: HRI-centered collaborative search.
% Primary comparison: B2 Reactive-Current vs B3 Predictive-Future.
% Primary system estimand: Delta_RMST at the common mission budget.
% This file is intentionally complete at the method/protocol level.
% It is designed to guide final platform freezing and the subsequent
% participant study. Numerical values that must come from frozen calibration
% or participant data should be inserted only after the corresponding
% pre-registered analysis has been completed.
% ============================================================

\newcommand{\figplaceholder}[2]{%
\fbox{%
\parbox[c][#1][c]{0.93\columnwidth}{%
\centering\footnotesize\textbf{FIGURE PLACEHOLDER}\\[2mm]#2}}}

\newcommand{\figplaceholderwide}[2]{%
\fbox{%
\parbox[c][#1][c]{0.94\textwidth}{%
\centering\footnotesize\textbf{FIGURE PLACEHOLDER}\\[2mm]#2}}}

\newcommand{\TBD}[1]{\textit{[#1]}}

% Figure workflow:
% 1) Draw the conceptual/method figures now and export them as vector PDF.
% 2) Put them under figures/ with the suggested names in this source.
% 3) Replace only the corresponding placeholder box with \includegraphics.
% 4) Leave result-figure comments untouched until formal experiments and
%    pre-registered analyses are complete.

\begin{document}

% ============================================================
% TITLE
% ============================================================
\title{Predictive Human Spatial Responsibility Redistribution for\\
Collaborative Multi-USV Probabilistic Search}

\author{Anonymous Authors}

\maketitle

% ============================================================
% ABSTRACT
% ============================================================
\begin{abstract}
This paper studies human--multi-USV collaboration in probabilistic maritime
search. The central human--robot interaction question is whether a robot fleet
should merely react to where a human-controlled vehicle is now, or anticipate
where that human intervention will allocate search effort in the near future.
We represent this effect as \emph{human spatial responsibility}: target-sensor
observations update a normalized Bayesian belief, whereas joystick, gaze, and
voice influence fleet planning without directly modifying the target
posterior. Under direct human control, a reactive-current baseline fixes the
human-controlled USV's current position $p_h(t)$ as a responsibility center;
the proposed predictive condition instead fixes the predicted future position
$\hat p_h(t+H)$. All searchable free cells remain in the same belief-weighted
allocation problem while only autonomous centers are updated. Pairwise
obstacle-aware A* paths are computed before path-cost Hungarian assignment.
Voice provides a separate explicit vehicle--destination responsibility channel,
while gaze is treated as optional predictive context rather than target
evidence. The evaluation freezes perception, sensor calibration, gaze
validation, and future-position prediction before mission testing. The primary
system outcome is the B2--B3 difference in restricted mean
time-to-confirmed-detection under a common mission budget; redundant sensing
overlap, success probability, fleet effort, NASA-TLX workload, and perceived
coordination alignment are secondary outcomes. An oracle future-position
condition estimates the performance gap due to prediction error independently
of the redistribution mechanism.
\end{abstract}

\begin{IEEEkeywords}
human--robot interaction, human--robot collaboration, multi-robot systems,
unmanned surface vehicles, probabilistic search, future-position prediction,
task allocation, maritime robotics
\end{IEEEkeywords}

% ============================================================
% I. INTRODUCTION
% ============================================================
\section{Introduction}

Multi-robot maritime search combines uncertain sensing with repeated spatial
allocation. Even when a coarse obstacle map is known before deployment, the
target location is unknown and onboard observations are imperfect. Bayesian
search therefore provides a probabilistically meaningful task state: target
belief is reduced only by calibrated positive and negative evidence rather
than by geometric visitation alone~\cite{b1,b7,b8}.

The central problem of this paper, however, is not Bayesian filtering itself.
It is how an autonomous fleet should adapt when a human becomes an active
member of the search team. A human may directly control one USV, explicitly
assign a vehicle to a destination, or provide spatial context through gaze.
These interventions change who is expected to search where. If the autonomous
fleet only reacts to the human-controlled USV's current position, it can still
allocate effort toward a region that the human is already moving toward. A
human-aware fleet should therefore reason about the future spatial consequence
of the intervention, not merely obey the intervention locally.

We call this consequence \emph{human spatial responsibility}: the future
search effort expected to be undertaken by the human-controlled or
human-commanded member of the fleet. Human spatial responsibility is not a
probability that the target is present. The target posterior and the human
responsibility representation answer different questions---\emph{where is the
target likely to be?} and \emph{where is the human likely to search?}---and
are deliberately kept separate until fleet-level allocation.

This leads to the primary research question:

\begin{quote}
\emph{Does predicted future human spatial responsibility provide actionable
coordination value beyond the current spatial responsibility of a
human-controlled USV?}
\end{quote}

The principal comparison is therefore deliberately narrow. Under direct human
control, the reactive-current condition represents human responsibility by the
current vehicle position $p_h(t)$. The predictive condition uses a frozen
future-position model to estimate $\hat p_h(t+H)$ and supplies that point to the
same fixed-center redistribution operator. An oracle condition uses the actual
future human-controlled position $p_h(t+H)$ offline to estimate the upper bound
of the same mechanism. Thus the scientific difference between the key
conditions is future human information, not a different planner.

Human behavior and intent prediction are established HRI topics~\cite{b2,b3,b4},
and Bayesian search, weighted clustering, A*, Hungarian assignment, speech
recognition, eye tracking, and Random Forest regression are established
components. The contribution is consequently not an algorithmic stack or a
multimodal interface. Instead, the paper makes three contributions:

\begin{enumerate}
\item a probabilistically consistent HRI formulation in which human
      intervention is represented as spatial search responsibility rather
      than target evidence;
\item a predictive fixed-center redistribution mechanism that uses the future
      spatial consequence of direct human control to reorganize the remaining
      autonomous fleet while retaining all searchable free cells in one
      belief-weighted allocation problem; and
\item a present--predicted--oracle evaluation ladder that compares current
      human responsibility, predicted future responsibility, and perfect
      future responsibility under the same redistribution operator, enabling
      prediction error to be separated from mechanism limitations.
\end{enumerate}

The primary hypothesis is mission-level HRI benefit: predictive-future
redistribution should reduce restricted mean time-to-confirmed-detection
relative to reactive-current redistribution. A secondary mechanism hypothesis
is that this improvement is mediated by reduced redundant sensing overlap
between the human-controlled and autonomous portions of the fleet. Human
workload and perceived coordination alignment are evaluated as secondary HRI
outcomes rather than used to redefine the primary system claim.

\begin{figure*}[t]
\centering
\figplaceholderwide{42mm}{
\textbf{HRI-centered system architecture.}\\
Target branch: RGB sensing $\rightarrow$ frozen detector $\rightarrow$
calibrated likelihood $\rightarrow$ Bayesian target belief.\\
Human branch: direct joystick control $+$ gaze context $\rightarrow$ future
position prediction $\hat p_h(t+H)$; voice $\rightarrow$ explicit
vehicle--destination responsibility.\\
Show three human responsibility representations: current $p_h(t)$, predicted
future $\hat p_h(t+H)$, and explicit commanded destination. The target and
human branches meet only at belief-aware fixed-center redistribution, followed
by pairwise A*, path-cost Hungarian assignment, execution, sensing, and
event-driven replanning. Add an explicit crossed arrow showing that human
signals do not directly update the target posterior.}
\caption{HRI-centered system concept. Calibrated target observations update the
Bayesian target belief, whereas direct control, gaze context, and explicit
voice commands determine current, predicted, or commanded human spatial
responsibility. The two information streams meet only at fleet allocation.}
\label{fig:architecture}
\end{figure*}

% ============================================================
% II. RELATED WORK
% ============================================================
\section{Related Work and Research Gap}

\subsection{Probabilistic Robotic Search}

Probabilistic search represents uncertainty about target location rather than
reducing search to geometric coverage. Bourgault \emph{et al.} presented
coordinated Bayesian search for a lost target~\cite{b1}. Dames showed how
finite field of view, false negatives, false positives, and probabilistic
search objectives can be integrated in multi-robot search and tracking
\cite{b7}. Jovan \emph{et al.} further demonstrated the importance of modeling
unreliable sensing when probabilistic estimates drive exploration~\cite{b8}.
In the present work, the Bayesian target posterior is not itself a novelty; it
provides the shared probabilistic state on which human-aware responsibility
redistribution operates.

A second distinction is important. A normalized posterior and accumulated
missed-detection evidence are not interchangeable. The former expresses
relative target plausibility after normalization, whereas the latter records
how compatible each cell remains with the accumulated negative observations.
Keeping both quantities avoids conflating ``visited'' with ``cleared.''

\subsection{Human Behavior Prediction and Anticipatory Coordination}

Intent inference and action prediction are established topics in human--robot
collaboration~\cite{b4}. Xie \emph{et al.} used inferred human search behavior
to support complementary human--robot search~\cite{b2}. Heintzman \emph{et
al.} incorporated anticipated human-searcher trajectories into search-and-
rescue planning~\cite{b3}. These studies define an important novelty boundary:
this work does not claim the first prediction of human search behavior.

The narrower gap addressed here is how the future spatial consequence of
direct human control should alter a belief-driven multi-robot responsibility
allocation. The present work does not claim task-level intent recognition:
the predictor estimates a future vehicle position. That position is interpreted
as a prospective human search responsibility and inserted as a fixed center
while all free cells remain in the same belief-weighted optimization.

\subsection{Collaborative Navigation and Task Allocation}

Human--robot collaborative navigation combines shared planning, communication,
authority changes, and task allocation. Dalmasso \emph{et al.} studied shared
task representation for collaborative navigation and search~\cite{b5}. Our
setting adds a dynamic fleet-allocation problem in which the same USV may be
autonomous, directly human controlled, or locked to an explicit human-issued
goal while the remaining vehicles continue probabilistic search.

Search responsibility is not treated as a navigation obstacle. A region
expected to be searched by the human-controlled USV remains traversable by
other USVs. Responsibility allocation and collision-free motion planning are
separate layers.

\subsection{Novelty Boundary}

This paper does not claim novelty for Bayesian filtering, Random Forest
regression, weighted Lloyd updates, A*, Hungarian assignment, speech
recognition, eye tracking, joystick teleoperation, ROS~2, Gazebo, or VRX
individually. The scientific hypothesis concerns the mechanism: predict the
future spatial consequence of human intervention, encode that future point as
a fixed search-responsibility center, and redistribute the remaining
autonomous fleet before a reactive condition receives the same spatial
responsibility explicitly.

% ============================================================
% III. PROBLEM FORMULATION
% ============================================================
\section{Problem Formulation}

\subsection{Environment, Target, and Fleet}

Let $\Omega\subset\mathbb{R}^2$ denote a bounded maritime workspace and
$\mathcal{O}\subset\Omega$ the known static obstacle region obtained from
preliminary reconnaissance. The searchable free space is

\begin{equation}
\mathcal{F}=\Omega\setminus\mathcal{O}.
\label{eq:free_space}
\end{equation}

The free space is discretized into $N$ searchable cells with centers

\begin{equation}
\mathcal{X}=\{x_1,\ldots,x_N\},\qquad x_i\in\mathcal{F}.
\label{eq:cells}
\end{equation}

Exactly one stationary target occupies one unknown searchable cell. Simulator
ground truth is retained for calibration and offline evaluation only; it is not
provided to the production detector, Bayesian estimator, clustering algorithm,
planner, or mission-success logic.

The fleet contains $R$ homogeneous USVs. Vehicle $r$ has planar position
$p_r(t)\in\mathbb{R}^2$ and heading $\psi_r(t)$. Its task-level authority state
is

\begin{equation}
a_r(t)\in\{\mathsf{A},\mathsf{H},\mathsf{L}\},
\label{eq:authority}
\end{equation}

where $\mathsf{A}$ denotes autonomous operation, $\mathsf{H}$ direct human
control, and $\mathsf{L}$ an explicit human-commanded locked task. Only
vehicles in $\mathsf{A}$ participate in the autonomous vehicle--center
assignment.

\subsection{Probabilistic Target State}

Let the discrete random variable $X$ denote the unknown target cell. Given
valid target-sensor observations $Z_{1:t}$, the normalized target belief is

\begin{equation}
b_i^t=P(X=x_i\mid Z_{1:t}),\qquad
\sum_{i=1}^{N}b_i^t=1.
\label{eq:posterior}
\end{equation}

Obstacle cells carry zero probability, and the prior $b_i^0$ is frozen at
mission start. Let $\bar D_{0:t}$ denote the event that no effective target
detection has occurred up to time $t$. The conditional cumulative
missed-detection probability is

\begin{equation}
S_i(t)=P(\bar D_{0:t}\mid X=x_i),\qquad S_i(0)=1.
\label{eq:cell_miss}
\end{equation}

The values $S_i(t)$ are not normalized across cells. The mission-level residual
miss probability is

\begin{equation}
R_{\mathrm{miss}}(t)=\sum_{i=1}^{N}b_i^0S_i(t).
\label{eq:rmiss}
\end{equation}

For histories containing only negative target evidence,

\begin{equation}
b_i^t=\frac{b_i^0S_i(t)}{R_{\mathrm{miss}}(t)}.
\label{eq:posterior_survival_relation}
\end{equation}

Equation~\eqref{eq:posterior_survival_relation} clarifies the difference
between absolute compatibility with accumulated negative evidence and relative
posterior probability. An unobserved cell can retain $S_i(t)=1$ while its
normalized posterior increases because mass is reduced elsewhere.

\subsection{Human Information Boundary}

Let $U_{1:t}^{h}$ denote human-interaction information, including joystick
actions, gaze behavior, interface-view context, semantic voice commands, and
predicted future human motion. The production target posterior in
Eq.~\eqref{eq:posterior} is defined only from calibrated target-sensor
evidence. Human interaction is not modeled as a target observation and
therefore does not directly modify $b_i^t$.

This restriction is deliberate. Human-derived information could enter the
Bayesian target state only after introducing and validating an explicit human-
observation likelihood. In the present framework, human information affects
search responsibility, not target probability.

\subsection{Human Spatial Responsibility and Search Allocation}

Human spatial responsibility denotes the portion of future search effort that
is expected to be undertaken by a human-controlled or human-commanded member of
the fleet. In the present point-responsibility formulation it is represented by
a fixed search center rather than a hard region deletion. Let $\mathcal{C}(t)$
denote the current set of search-responsibility centers.
Given the target posterior, the spatial allocation objective is

\begin{equation}
\Phi(\mathcal{C};t)=
\sum_{i=1}^{N}
b_i^t\min_{c\in\mathcal{C}}\|x_i-c\|_2^2.
\label{eq:responsibility_objective}
\end{equation}

Human-related fixed centers are collected in $\mathcal{C}_F$ and autonomous
centers in $\mathcal{C}_A$. If $K_A$ autonomous vehicles are available,

\begin{equation}
\begin{aligned}
\mathcal{C}_A^{\star}
=\arg\min_{\mathcal{C}_A}
&\;\Phi(\mathcal{C}_F\cup\mathcal{C}_A;t)\\[-1mm]
\text{s.t. }&\mathcal{C}_A\subseteq\mathcal{F},\quad
|\mathcal{C}_A|=K_A.
\end{aligned}
\label{eq:constrained_responsibility}
\end{equation}

All searchable free cells remain in the optimization. Fixed centers do not
move; only autonomous centers are updated. No residual human region is manually
removed.

\subsection{Reactive-Current and Predictive-Future Collaboration}

The principal HRI study considers one directly human-controlled USV at a time
and the remaining autonomous fleet. Let $p_h(t)$ denote the current position of
the human-controlled USV. The reactive-current condition represents the
human's present spatial responsibility by

\begin{equation}
c_h^{B2}(t)=p_h(t).
\label{eq:b2_current_center}
\end{equation}

The proposed predictive-future condition estimates the same vehicle's future
position at horizon $H$ and uses

\begin{equation}
c_h^{B3}(t)=\hat p_h(t+H).
\label{eq:b3_predicted_center}
\end{equation}

For offline oracle evaluation, perfect future responsibility is represented by

\begin{equation}
c_h^{B4}(t)=p_h(t+H).
\label{eq:b4_oracle_center}
\end{equation}

B2, B3, and B4 use the same constrained redistribution operator in
Eq.~\eqref{eq:constrained_responsibility}. Their scientific difference is the
human spatial information available to the planner: present position,
predicted future position, or actual future position. The prediction neither
commands the human-controlled USV nor claims to infer the operator's final
semantic goal.

Explicit voice commands form a separate responsibility channel. A valid
vehicle--destination command produces a locked pair and an explicit fixed
center; it is not used to define the primary B2--B3 comparison and is not
provided as a privileged input to the future-position predictor.

% ============================================================
% IV. PROBABILISTIC TARGET SEARCH
% ============================================================
\section{Probabilistic Target Search}

\subsection{Truth-Free Visual Observation Interface}

The formal search system uses a frozen single-class visual target detector. The
detector architecture is an implementation choice; the search algorithm
consumes only the detector outputs and their independently calibrated
statistical model. Simulator target truth is used for detector development,
calibration, and offline evaluation, but it does not enter production
inference.


% ------------------------------------------------------------
% FIG. 2: SENSOR--BAYES CLOSED LOOP
% Suggested final file: figures/fig2_sensor_bayes_pipeline.pdf
% Replace only the placeholder box when the final vector figure is ready.
% ------------------------------------------------------------
\begin{figure*}[t]
\centering
\figplaceholderwide{46mm}{
\textbf{Visual observation $\rightarrow$ Bayesian search closed loop.}\\
Online branch: RGB image $\rightarrow$ frozen detector $\rightarrow$
positive/negative observation. Positive branch: bbox/pixel anchor
$\rightarrow$ image-time camera ray $\rightarrow$ water-plane localization
$\rightarrow z_k$. Negative branch: range/FOV + static LOS
$\rightarrow$ missed-detection likelihood. Both branches enter
$L_i(Z_k)\rightarrow b^{k+1}\rightarrow$ belief-aware centers
$\rightarrow A^*\rightarrow$ path-cost Hungarian $\rightarrow$ USV motion
$\rightarrow$ next observation.\\
Offline branch: independent calibration of
$P_D$, $P_{\rm FA}$, $q_T$, $\nu_{\rm eff}$, and $\Gamma$, with a dashed
arrow into the online likelihood model.}
\caption{Closed-loop visual probabilistic search pipeline. The visual detector
does not directly modify the target posterior. Its outputs are interpreted
through an independently calibrated observation model, and the resulting
posterior drives the subsequent fleet-planning cycle.}
\label{fig:sensor_bayes}
\end{figure*}

As summarized in Fig.~\ref{fig:sensor_bayes}, the perception frontend,
probabilistic observation model, Bayesian estimator, and fleet planner form a
closed loop but remain distinct modules. In particular, the detector score is
not interpreted as target probability, and offline simulator truth is not an
online inference input.

Let $C_r$ denote the optical center of camera $r$. For candidate cell $x_i$,
define the horizontal camera-relative geometry

\begin{equation}
\begin{aligned}
\rho_{ri}(t)&=\|x_i-p_{C_r}(t)\|_2,\\
\beta_{ri}(t)&=
\operatorname{wrap}\!\left[
\operatorname{atan2}(\Delta y,\Delta x)-\psi_{C_r}(t)
\right],
\end{aligned}
\label{eq:camera_geometry}
\end{equation}

where $\Delta=[\Delta x,\Delta y]^\top=x_i-p_{C_r}(t)$ and $\psi_{C_r}$ is the
camera optical heading projected into the horizontal plane.


The geometry $(\rho_{ri},\beta_{ri})$ does not require a stereo or depth
camera. It is evaluated for each \emph{hypothesized} target cell $x_i$ from
the known map coordinate of that cell and the time-aligned camera pose. Depth
is required only when a positive image candidate is converted to a world
measurement; that localization uses the ray--water-plane construction
described below.

\subsection{Static Visibility Qualification}

Static obstacles can invalidate negative visual evidence. Let

\begin{equation}
V_{ri}(t)\in\{\mathsf{VIS},\mathsf{OCC},\mathsf{UNK}\}
\label{eq:visibility_state}
\end{equation}

represent the static visibility state between camera $r$ and candidate cell
$i$. The state is computed from verified physical obstacle geometry, not from
the safety-inflated navigation map. Static line of sight is qualified using the
known environment model; target ground truth is not an input.

For $V_{ri}=\mathsf{VIS}$, the calibrated detector supplies an end-to-end
probability of an effective correct target observation,

\begin{equation}
P_{D,ri}=P_D(\rho_{ri},\beta_{ri}).
\label{eq:pd_generic}
\end{equation}

The function $P_D(\rho,\beta)$ is identified from independent calibration data.
A logistic or other parametric form is retained only if supported by the
calibration fit; the paper does not treat historical engineering parameters as
scientific constants.

When $V_{ri}=\mathsf{OCC}$ or $\mathsf{UNK}$, the current observation is not
qualified to reduce the probability of cell $i$. This is an abstention rule,
not an assertion that the physical detector has zero detection probability.
Dynamic vessel occlusion and transient wave-crest occlusion are treated as
residual sensor uncertainty unless explicitly modeled in calibration.

\subsection{Negative Observation Update}

Let $\nu_{\rm eff}$ denote the calibrated effective rate of statistically
independent observation opportunities. For approximately constant geometry
over interval $\Delta t$, define

\begin{equation}
\hspace{-0.8em}
\begin{aligned}
m_{ri}(t,\Delta t)
&=
\begin{cases}
[1-P_{D,ri}(t)]^{\nu_{\rm eff}\Delta t},
V_{ri}=\mathsf{VIS},\\
1,\hspace{40px}
V_{ri}\in\{\mathsf{OCC},\mathsf{UNK}\}.
\end{cases}
\end{aligned}
\label{eq:interval_miss}
\end{equation}

The corresponding first-effective-detection hazard for visible cells is

\begin{equation}
h_{ri}(t)=-\nu_{\rm eff}\ln[1-P_{D,ri}(t)],
\label{eq:detection_hazard}
\end{equation}

so that $m_{ri}=\exp[-h_{ri}\Delta t]$ under constant geometry.

Assuming conditional independence across USVs given the target cell, the
fleet-level missed-detection likelihood is

\begin{equation}
M_i(t)=\prod_{r\in\mathcal{R}}m_{ri}(t,\Delta t).
\label{eq:fleet_miss}
\end{equation}

For an interval with no valid positive target observation,

\begin{equation}
\begin{aligned}
b_i^{t+1}
&=\frac{b_i^tM_i(t)}{\sum_j b_j^tM_j(t)},\\
S_i(t+1)&=S_i(t)M_i(t).
\end{aligned}
\label{eq:negative_update}
\end{equation}

If the per-opportunity false-alarm probability $P_{\rm FA}$ is independent of
candidate cell, its common factor cancels in the normalized negative update.
Thus missed detections reduce belief only where the target should have been
observable.

\subsection{World Localization of Positive Candidates}

A detector bounding box is converted to a world-frame measurement through a
camera model and the nominal water surface. Let $K$ be the camera intrinsic
matrix and $(u,v)$ the selected image anchor. The camera-optical ray is

\begin{equation}
d_{C_{\rm opt}}
=\frac{K^{-1}[u,v,1]^\top}
{\|K^{-1}[u,v,1]^\top\|_2}.
\label{eq:camera_ray}
\end{equation}

The image/CameraInfo frame is a generic camera-sensor frame rather than the ROS
optical frame. The global optical transform is therefore

\begin{equation}
T^G_{C_{\rm opt}}(t_k)=
T^G_B(t_k)
T^B_{C_{\rm sensor}}(t_k)
T^{C_{\rm sensor}}_{C_{\rm opt}},
\label{eq:camera_transform_chain}
\end{equation}

where all time-varying terms are evaluated at the image measurement time.
With global camera origin $o_G$ and ray direction $d_G$, the nominal water
plane is

\begin{equation}
n^\top p+d=0.
\label{eq:water_plane}
\end{equation}

The valid intersection is

\begin{equation}
\lambda=-\frac{n^\top o_G+d}{n^\top d_G},\qquad
z_k=o_G+\lambda d_G,
\label{eq:ray_plane}
\end{equation}

provided the ray is not parallel to the plane and $\lambda>0$. Anchor choice is
frozen from independent localization residual validation rather than chosen for
visual convenience.

\subsection{Positive Observation Update}

Let $z_k$ denote a geometrically valid world-frame positive candidate. The
formal positive likelihood is

\begin{equation}
L_i^+(z_k)=P(z_k,\mathrm{positive}\mid X=x_i),
\label{eq:positive_likelihood_interface}
\end{equation}

whose target-generated localization component is denoted
$q_T(z_k\mid X=x_i)$. The detector operating point, false-alarm behavior, and
localization residual model are estimated from independent calibration data.
The Bayesian update is

\begin{equation}
b_i^+=
\frac{b_i^-L_i^+(z_k)}{\sum_j b_j^-L_j^+(z_k)}.
\label{eq:positive_update}
\end{equation}

A detector confidence score is not treated as a target posterior probability.
If the perception frontend retains more than one simultaneous candidate, the
positive likelihood must use a set-valued clutter model rather than reusing a
single-candidate approximation.

\subsection{Unified Observation Model}

Positive and negative observations are two outcomes of one Bayesian sensor
model. For effective observation $Y_k$,

\begin{equation}
\hspace{-0.8em}
\begin{aligned}
b_i^{k+1}
&=\frac{b_i^kL_i(Y_k)}{\sum_j b_j^kL_j(Y_k)},\\[-0.3ex]
L_i(Y_k)
&=P(Y_k\mid X=x_i).
\end{aligned}
\label{eq:unified_bayes}
\end{equation}

Negative observations use the visibility-qualified missed-detection
likelihood. Positive observations use the calibrated localized likelihood in
Eq.~\eqref{eq:positive_likelihood_interface}.

\subsection{Confirmation and Episode Termination}

Detector output, Bayesian belief, and mission success are distinct quantities.
Let $\tau_{\rm det}$ denote the frozen detector operating threshold and let
$\Gamma_k\in\{0,1\}$ denote the independently calibrated confirmation rule
applied to eligible candidate history. The formal confirmed-detection time is

\begin{equation}
T_{\rm conf}=\inf\{t_k:\Gamma_k=1\}.
\label{eq:tconf}
\end{equation}

The primary experiment ends at

\begin{equation}
T_{\rm end}=\min\{T_{\rm conf},T_{\rm budget}\}.
\label{eq:tend}
\end{equation}

A run is successful only if calibrated confirmation occurs within budget:

\begin{equation}
\delta_{\rm succ}=
\mathbf{1}[T_{\rm conf}\leq T_{\rm budget}].
\label{eq:success_indicator}
\end{equation}

\texttt{BUDGET\_EXPIRED}, \texttt{SEARCH\_EXHAUSTED}, \texttt{STOPPED}, and
\texttt{INVALID} are non-success terminal causes. Ground-truth proximity,
posterior-MAP proximity, and $R_{\rm miss}$ thresholds remain offline diagnostics
and are not formal target-detection events.

% ============================================================
% V. MULTIMODAL HUMAN SPATIAL RESPONSIBILITY
% ============================================================
\section{Human Spatial Responsibility and Interaction Channels}

\subsection{Interaction Semantics}

The interaction modalities support different HRI roles and are not treated as
an undifferentiated multimodal-fusion contribution. Joystick input directly
controls the selected human-controlled USV and provides implicit motor evidence
about its future spatial responsibility. Gaze supplies attentional and spatial
context that may improve future-position prediction when its task-space
mapping and timing are valid. Voice supplies explicit spatial task
responsibility in the restricted form of a vehicle--destination command.

None of these signals is target evidence. Joystick and gaze may enter the
future-position predictor; an explicit voice destination does not, because the
responsibility is already specified. A valid voice command creates a locked
vehicle--destination pair, excludes that USV from autonomous Hungarian
assignment, and inserts the commanded destination as a fixed responsibility
center for the remaining fleet.

% ------------------------------------------------------------
% FIG. 3: HUMAN SPATIAL RESPONSIBILITY CHANNELS
% Suggested final file: figures/fig3_human_responsibility.pdf
% ------------------------------------------------------------
\begin{figure}[t]
\centering
\figplaceholder{45mm}{
\textbf{Human spatial responsibility channels.}\\
Joystick $\rightarrow$ direct control / implicit motor intent $\rightarrow$
future-position predictor;\\
Gaze $\rightarrow$ attentional/spatial context $\rightarrow$ optional predictor
features;\\
Voice $\rightarrow$ explicit vehicle--destination responsibility $\rightarrow$
locked fixed center.\\
Show the predictor output $\hat p_h(t+H)$ as predicted responsibility and the
current position $p_h(t)$ as reactive-current responsibility. Add an explicit
``does not update target belief'' boundary.}
\caption{Human interaction channels and their HRI roles. Direct control and
valid gaze context support implicit future-responsibility prediction, whereas
voice provides explicit vehicle--destination responsibility. Human information
constrains planning but does not directly update the Bayesian target belief.}
\label{fig:human_responsibility}
\end{figure}

\subsection{Tobii Eye-Tracking Pipeline}

The eye tracker produces screen-space gaze, not world-space target evidence.
Let the measured gaze be

\begin{equation}
g_s(t)=[u_g(t),v_g(t)]^\top,
\label{eq:screen_gaze}
\end{equation}

with validity and source timestamps. The acquisition path preserves tracker
time, Windows-side sequence information, Linux receive time, ROS publication
time, active interface view, and the currently viewed and controlled vehicles.
These timestamps are aligned to the experiment time base before causal feature
construction.

Two calibrations are distinguished. Device calibration maps eye measurements
to screen coordinates and is performed with the eye-tracker manufacturer's
procedure. Task-space validation then checks the experiment-specific mapping
from screen gaze to interface coordinates and, in the global map view, to world
coordinates. Participant sessions begin only after both checks satisfy the
pre-registered accuracy and validity criteria.

\subsection{Task-Space Gaze Mapping}

In the global overview, valid screen gaze is mapped through the actual Qt map
rectangle to a world point $p_g^W(t_g)$. Relative to the human-controlled USV,

\begin{equation}
\Delta p_g^W(t_g)=p_g^W(t_g)-p_h^W(t_g),
\label{eq:gaze_relative_world}
\end{equation}

and the body-frame gaze vector is

\begin{equation}
g_G^B(t_g)=R[-\psi_h(t_g)]\Delta p_g^W(t_g),
\label{eq:gaze_body}
\end{equation}

where

\begin{equation}
R(\psi)=
\begin{bmatrix}
\cos\psi&-\sin\psi\\
\sin\psi&\cos\psi
\end{bmatrix}.
\label{eq:rotation}
\end{equation}

Gaze outside the displayed map rectangle remains invalid rather than being
clamped to the nearest map point.

In FPV mode, gaze is retained in image coordinates. A normalized horizontal
offset is

\begin{equation}
\tilde u_g=\frac{u_g-u_c}{W_{\rm img}},
\label{eq:fpv_gaze}
\end{equation}

with an analogous normalized vertical offset when used. FPV gaze is not
silently projected to world $x$--$y$ unless a validated camera-ray mapping is
explicitly invoked. The active view mode is observed context, not a target
label.

\subsection{Gaze Validation}

Task-space validation presents known screen or map targets $q_k$ and records the
corresponding estimated gaze locations $\hat q_k$. Error is

\begin{equation}
e_k=\|\hat q_k-q_k\|_2.
\label{eq:gaze_error}
\end{equation}

For global-map targets, the same validation is expressed in world meters.
Median error, upper-tail error, valid-sample fraction, and out-of-map fraction
are reported. Cross-host clock mapping is validated separately by measuring the
residual timing error after synchronization; gaze samples that cannot be
aligned within the frozen tolerance are excluded from predictor windows rather
than assigned an arbitrary time.

% ============================================================
% VI. HUMAN FUTURE-POSITION PREDICTION
% ============================================================
\section{Human Future-Position Prediction}

The predictor estimates the future spatial consequence of direct human
control. It is therefore a future-position predictor, not a classifier of
high-level task semantics such as search, hold, patrol, or pursuit.

\subsection{Prediction Target}

For causal history duration $W$ and horizon $H$, the future displacement of the
human-controlled USV is expressed in the vehicle body frame at time $t$,

\begin{equation}
\Delta p_h^B(t,H)=
R[-\psi_h(t)]\left[p_h(t+H)-p_h(t)\right].
\label{eq:prediction_target}
\end{equation}

The body-frame representation reduces dependence on absolute map orientation.
Only samples contained within a single human-controlled segment are valid;
windows crossing authority changes, controlled-vehicle changes, mission resets,
or simulation resets are rejected.

\subsection{Causal Feature Construction}

The recorder stores an asynchronous research-data superset. Candidate feature
family $\mathcal{F}_k$ maps the interval $[t-W,t]$ into

\begin{equation}
x_t^{(k)}=\mathcal{F}_k\!\left(\mathcal{D}_{\rm raw}[t-W,t]\right),
\qquad x_t^{(k)}\in\mathbb{R}^{d_k}.
\label{eq:feature_mapping}
\end{equation}

The stable control/dynamics baseline summarizes body velocity, speed, yaw rate,
body accelerations, semantic manual commands, and left/right thrust using last,
mean, standard deviation, and slope statistics. Multimodal candidate families
add gaze, view context, and switch/authority context. Voice-specified explicit
goals are not included as privileged predictor inputs because they already
constitute explicit task intent.

\begin{table}[t]
\centering
\caption{Predefined predictor feature families.}
\label{tab:feature_families}
\small
\setlength{\tabcolsep}{4pt}
\begin{tabular}{clc}
\toprule
ID & Information & Dim.\\
\midrule
$F_0$ & Reduced control/dynamics & 20\\
$F_1$ & Full control/dynamics & 40\\
$F_2$ & $F_1$ + gaze & 54\\
$F_3$ & $F_2$ + view context & 58\\
$F_4$ & $F_3$ + switch context & 61\\
\bottomrule
\end{tabular}
\end{table}

The final feature family is selected using participant-disjoint validation and
then frozen before mission testing. Thus the table defines a predeclared
ablation family rather than asserting that the largest representation is best.

\subsection{Random-Forest Regression}

The predictor estimates body-frame future displacement,

\begin{equation}
\widehat{\Delta p}_h^B(t,H)=f_{\rm RF}(x_t^{\star}),
\label{eq:rf_regression}
\end{equation}

where $x_t^{\star}$ denotes the frozen selected feature family. The world-frame
future point is

\begin{equation}
\hat p_h(t+H)=
p_h(t)+R[\psi_h(t)]\widehat{\Delta p}_h^B(t,H).
\label{eq:world_prediction}
\end{equation}

Random Forest is used as an interpretable regression component rather than
claimed as the main contribution. Its role is to provide a reproducible future
spatial consequence estimate for the planning experiment.

\subsection{Prediction Quality and Coordination Lead Time}

Final displacement error is

\begin{equation}
\mathrm{FDE}(H)=
\|\hat p_h(t+H)-p_h(t+H)\|_2.
\label{eq:fde}
\end{equation}

Median, 90th, and 95th percentile FDE are reported because occasional large
errors can trigger inappropriate redistribution. Let $T_{\rm sys}$ denote the
measured delay from prediction availability to commitment of the corresponding
fleet plan. Effective anticipation time is

\begin{equation}
H_{\rm eff}=H-T_{\rm sys}.
\label{eq:heff}
\end{equation}

A horizon with $H_{\rm eff}\le0$ cannot provide proactive coordination even if
its FDE is small. The final $(W,H,F_s)$ is therefore selected using held-out
participants by jointly considering prediction error, valid-window support,
and positive coordination lead time.

% ============================================================
% VII. PREDICTIVE SPATIAL RESPONSIBILITY REDISTRIBUTION
% ============================================================
\section{Predictive Spatial Responsibility Redistribution}

\subsection{Fixed-Center Belief-Aware Update}

Equation~\eqref{eq:constrained_responsibility} is implemented using a
weighted Lloyd-style update with fixed centers. Each free cell is assigned to
its nearest current center. For autonomous center $c_j$, let $\mathcal{V}_j$
denote its assigned cells. The weighted centroid is

\begin{equation}
\bar c_j=
\frac{\sum_{i\in\mathcal{V}_j}b_i^t x_i}
{\sum_{i\in\mathcal{V}_j}b_i^t},
\label{eq:weighted_centroid}
\end{equation}

and the updated autonomous center is projected to an admissible free cell in
its cluster,

\begin{equation}
c_j^+=\Pi_{\mathcal{F}\cap\mathcal{V}_j}(\bar c_j).
\label{eq:center_projection}
\end{equation}

Fixed centers are never moved. Crucially, all non-obstacle free cells
participate again whenever reclustering occurs. The method does not remove a
previous human region from the optimization.

\subsection{Current, Predicted, and Explicit Human Responsibility}

During direct human control, the reactive-current condition inserts the current
human-controlled position $p_h(t)$ into $\mathcal{C}_F$, while the predictive
condition inserts $\hat p_h(t+H)$. The human-controlled USV is excluded from
autonomous matching in both conditions. This design ensures that B2 and B3
share the same responsibility-redistribution operator and differ only in
whether the planner uses present or predicted future human spatial
responsibility.

An explicit voice destination is handled separately. It is inserted directly
into $\mathcal{C}_F$, and the associated vehicle--destination pair remains
locked for the command interval. Explicit commands therefore validate the same
responsibility interface without defining the predictive hypothesis.

\subsection{Obstacle-Aware Path Planning and Assignment}

For each autonomous USV $r$ and autonomous center $c_j$,

\begin{equation}
\mathcal{P}_{rj}=A^*(p_r,c_j;\mathcal{M}).
\label{eq:astar_path}
\end{equation}

If $\mathcal{P}_{rj}=(q_0,\ldots,q_{n_{rj}-1})$, its path length is

\begin{equation}
L(\mathcal{P}_{rj})=
\sum_{\ell=0}^{n_{rj}-2}\|q_{\ell+1}-q_{\ell}\|_2.
\label{eq:path_length}
\end{equation}

The assignment cost is

\begin{equation}
J_{rj}=L(\mathcal{P}_{rj}),
\label{eq:path_cost}
\end{equation}

with infeasible pairs removed from the admissible matching set. The remaining
minimum-cost one-to-one assignment is

\begin{equation}
\pi^*=\arg\min_{\pi}
\sum_{r\in\mathcal{R}_A}J_{r,\pi(r)}.
\label{eq:assignment}
\end{equation}

All eligible pairwise A* paths are computed before Hungarian matching. The
selected path is the path subsequently executed by the controller, ensuring
consistency between assignment cost and execution.

\subsection{Event-Driven Replanning}

Replanning is triggered by task-relevant events rather than every incoming
message. Events include autonomous arrival, explicit human-goal changes,
predicted future-point changes, controlled-USV changes, command
completion/cancellation/replacement, invalid allocation, and invalid path
state. Belief-only replanning uses a center-set displacement gate to avoid
chattering when posterior changes are negligible.

Let $\hat p_h^{\rm plan}(t+H)$ denote the prediction associated with the active
plan and $\hat p_h^{\rm new}(t+H)$ the latest prediction. A prediction-triggered
replan occurs when

\begin{equation}
\left\|\hat p_h^{\rm new}(t+H)-
\hat p_h^{\rm plan}(t+H)\right\|_2
\geq\delta_{\rm pred}.
\label{eq:prediction_trigger}
\end{equation}

The threshold $\delta_{\rm pred}$ is frozen from held-out prediction-jitter and
prediction-error statistics rather than tuned on mission outcomes.

% ------------------------------------------------------------
% FIG. 4: CORE CONTRIBUTION
% Suggested final file: figures/fig4_predictive_redistribution.pdf
% Recommended layout: three map panels + one compact timing strip.
% ------------------------------------------------------------
\begin{figure*}[t]
\centering
\figplaceholderwide{50mm}{
\textbf{Core HRI mechanism: current versus predicted human responsibility.}\\
Panel (a): before takeover, autonomous belief-aware centers.\\
Panel (b): B2 Reactive-Current, with the human-controlled USV at $p_h(t)$ and
that current position fixed as the human responsibility center; remaining
autonomous centers are re-optimized.\\
Panel (c): B3 Predictive-Future, with $\hat p_h(t+H)$ fixed instead; show how
the remaining autonomous centers move earlier to reduce prospective overlap.\\
Optional inset: B4 oracle center $p_h(t+H)$ for offline upper-bound analysis.\\
Show that all free cells remain active in every panel. Bottom timing strip:
current time $t$, prediction horizon $t+H$, prediction availability, plan
commitment, and effective anticipation time $H_{\rm eff}$.}
\caption{Current versus predicted human spatial responsibility. B2, B3, and
the offline B4 oracle use the same fixed-center redistribution operator; they
differ only in whether the fixed human center is $p_h(t)$,
$\hat p_h(t+H)$, or $p_h(t+H)$.}
\label{fig:reactive_predictive}
\end{figure*}

% ============================================================
% VIII. SYSTEM IMPLEMENTATION
% ============================================================
\section{System Implementation}

\subsection{ROS~2, Gazebo, and VRX}

The platform is implemented in ROS~2 Humble, Gazebo, and VRX with ten
homogeneous WAM-Vs. A single launch entry point starts the world, overview
state, perception, search, planning, controller, interaction, and recording
layers. The research recorder is the unique authoritative experiment-data
writer; raw records remain immutable after finalization.

\subsection{Canonical Search Target and World Suite}

Formal experiments use a versioned stationary floating search target
(``search\_target'') rather than a scenario-specific object name. The benchmark
asset has fixed episode-level planar position and a geometry that intersects
the nominal water surface. Target asset, camera configuration, and world
version are hashed and stored in each experiment manifest.

Four worlds separate sensing identification from mission-level mechanism tests.
The calibration world contains open water and is used primarily for controlled
sensor characterization; it is not used as an HRI benchmark. The open-search world weakens obstacle effects and
provides a baseline probabilistic-search environment. The channel world
introduces path-distortion structures that make obstacle-aware assignment
relevant. The archipelago world contains several semi-open search regions and
is used as the principal environment for human spatial-responsibility
redistribution. These worlds represent different experimental mechanisms,
not a claimed monotonic difficulty scale.

The nominal search domain is an $800\times800$~m local ENU region represented
by a coarse $10$~m planning grid. Physical obstacle geometry, visual geometry,
LOS geometry, and navigation inflation are tracked separately. Static LOS uses
uninflated opaque physical polygons; the A* grid uses conservative inflation
for vehicle clearance.

\begin{table}[t]
\centering
\caption{Role of the frozen world suite.}
\label{tab:worlds}
\small
\setlength{\tabcolsep}{3.5pt}
\begin{tabular}{ll}
\toprule
World & Primary experimental role\\
\midrule
Calibration & $P_D$, $P_{FA}$, localization, temporal sensing\\
Open search & probabilistic-search baseline\\
Channel & obstacle/path-cost relevance\\
Archipelago & spatial-responsibility redistribution\\
\bottomrule
\end{tabular}
\end{table}

% ------------------------------------------------------------
% FIG. 5: WORLD SUITE
% Suggested final file: figures/fig5_world_suite.pdf
% Four aligned overhead panels: Calibration / Open / Channel / Archipelago.
% Show navigable water, physical obstacles, representative starts, and target
% allowed region. Do not encode results or imply a monotonic difficulty order.
% ------------------------------------------------------------
\begin{figure*}[t]
\centering
\figplaceholderwide{43mm}{
\textbf{Versioned maritime search environments.}\\
Four aligned top-down panels: (a) Calibration, (b) Open, (c) Channel,
(d) Archipelago. Use the same map scale and visual legend. Distinguish physical
obstacle geometry from navigation inflation if both are shown. The figure
should communicate each world's experimental role rather than an
easy-to-hard ordering.}
\caption{Representative benchmark environments used to separate sensor
characterization, low-obstacle probabilistic search, path-sensitive
allocation, and spatial-responsibility redistribution.}
\label{fig:world_suite}
\end{figure*}

% PLATFORM GATE BEFORE FORMAL MISSION STUDY:
% Freeze the learned detector checkpoint and runtime preprocessing, then perform
% independent sensor calibration. Do not insert engineering detector constants
% into the final paper as calibration results.
\subsection{Visual Perception and World Localization}

A single detector owner serves all camera streams. The formal perception
frontend is a frozen single-class learned detector whose checkpoint, input
preprocessing, operating point, and inference configuration are fixed before
sensor calibration. Detector development uses group-disjoint training,
validation, and detector-test data. Sensor calibration and mission data are
separate roles and cannot be used for detector training or threshold tuning.

World localization uses the image-time transform chain in
Eq.~\eqref{eq:camera_transform_chain}. The implementation explicitly separates
the generic Gazebo camera sensor frame from the ROS optical frame, avoiding a
frame-convention error in back-projected rays. Missing CameraInfo, unavailable
transforms, stale timing, parallel rays, and behind-camera intersections fail
closed and do not become negative target evidence.

\subsection{Joystick, Voice, and Authority}

Joystick takeover assigns one vehicle to direct human control. Raw joystick
input, semantic manual command, authority state, and executed thrust are
recorded separately. Voice is transcribed and normalized into a deliberately restricted canonical
command representation containing a vehicle and a destination. This channel is
used to validate explicit spatial responsibility, not general high-level
language reasoning. A valid vehicle--destination command creates a locked pair
when execution is enabled. Locked pairs use the same A* path planner but are
excluded from autonomous Hungarian assignment.


% PLATFORM GATE BEFORE LIVE HRI:
% Freeze a single world-independent semantic-destination resolver/rejection
% rule for command regions that intersect obstacles. Do not use map-specific
% parser branches or silently relocate commands without provenance.

% PLATFORM GATE BEFORE PARTICIPANT STUDY:
% Complete Tobii device calibration, task-space validation, and cross-host
% time-alignment validation. Raw UDP reception alone is not sufficient.
\subsection{Eye Tracking and Time Alignment}

Tobii gaze is acquired on the operator workstation and transmitted with source
validity, tracker timestamp, sequence metadata, and screen geometry. The ROS
bridge preserves source and receive timing; the operator interface supplies
view mode and screen-to-map geometry. Device calibration and task-space
validation are performed before experimental sessions. Cross-host timing is
mapped to the experiment time base and checked against a predeclared residual
alignment tolerance. Missing or invalid gaze remains missing rather than being
zero-filled.

\subsection{Research Data Architecture}

The data model distinguishes record order, measurement time, and source-local
order. Predictor windows are generated offline from immutable raw streams using
mission-relative simulation time plus validated source-time alignment where
needed. Windows never cross mission reset, authority change, controlled-vehicle
change, or human-controlled segment boundaries. This organization allows the
same raw participant sessions to support alternative causal history windows and
feature-family ablations without recollecting data.

% ============================================================
% IX. EXPERIMENTAL METHODOLOGY
% ============================================================
\section{Experimental Methodology}

The experimental protocol is structured so that perception development,
sensor-model identification, gaze validation, predictor selection, and
mission-level hypothesis testing cannot leak into one another. All quantities
used by the final participant study are frozen before the corresponding test
runs.

\subsection{Stage A: Detector Development and Freeze}

Detector development uses disjoint training and validation episode groups.
A separate detector-test set remains sealed during development. The detector
checkpoint, image preprocessing, class mapping, non-maximum suppression,
operating score range, and runtime adapter are frozen before Stage~B. Detector
training details are not a scientific contribution; the paper reports only the
frozen detector identity and held-out detector characterization required to
make the sensing model reproducible.

\subsection{Stage B: Sensor Calibration}

\subsubsection{Detection probability}

Target-present trials span controlled camera-relative range and bearing. For
each calibration opportunity, simulator truth is used offline to classify the
frozen detector output as TP or FN. Empirical detection probability is

\begin{equation}
\widehat P_D(\rho,\beta)=
\frac{N_{\rm TP}(\rho,\beta)}
{N_{\rm TP}(\rho,\beta)+N_{\rm FN}(\rho,\beta)}.
\label{eq:empirical_pd}
\end{equation}

A parametric model is fitted only if it passes goodness-of-fit and residual
checks. Confidence intervals or posterior uncertainty are reported for the
fitted parameters.

\subsubsection{False alarms and operating point}

Target-absent and distractor-present trials estimate false-positive behavior,

\begin{equation}
\widehat P_{\rm FA}=
\frac{N_{\rm FP}}{N_{\rm FP}+N_{\rm TN}}.
\label{eq:empirical_pfa}
\end{equation}

The final detector operating point is selected before mission testing to
satisfy a predeclared false-alarm or false-termination constraint while
retaining useful detection probability. Mission outcomes are not used to
retune this threshold.

\subsubsection{Localization and anchor selection}

For true-positive detections,

\begin{equation}
e_{z,k}=z_k-p_{{\rm GT},k}
\label{eq:localization_error}
\end{equation}

is used to compare admissible bounding-box anchors and to estimate localization
bias, covariance, and range/bearing dependence. The final anchor and
$q_T(z\mid X=x_i)$ are frozen from these residuals. A Gaussian model is used
only if supported by residual diagnostics.

\subsubsection{Temporal dependence}

Continuous detector sequences at fixed relative geometry are used to estimate
frame-to-frame dependence and $\nu_{\rm eff}$. First-detection-time survival
under the selected temporal model is validated directly rather than inferred
from camera frame rate.

\subsubsection{Confirmation rule}

The formal confirmation rule $\Gamma_k$ is selected from held-out calibration
sequences. A single validated detection is allowed only if its mission-level
false-termination risk satisfies the predeclared criterion; otherwise repeated
spatiotemporally consistent positives are required. The final rule is frozen
before any B0--B4 mission comparison.

\subsection{Stage C: Tobii and Interface Validation}

Each participant first completes the manufacturer's device calibration. A
second task-space validation presents known targets in the actual experiment
interface. For the global overview, validation points span the map display;
for FPV, validation checks image-space gaze. The study records median and
upper-tail task-space error, valid-sample fraction, out-of-map fraction, and
cross-host timing residual. Sessions that do not satisfy the pre-registered
quality criteria are recalibrated before data collection.

This stage is essential because valid eye-tracker UDP traffic alone does not
establish scientifically valid gaze features. The study validates the complete
measurement chain used by the predictor.

% PLATFORM GATE BEFORE B2--B3 TESTING:
% Use participant-disjoint data to select/freeze W, H, F_s, feature family,
% RF configuration, and delta_pred. Mission test data must remain untouched.
\subsection{Stage D: Predictor Development}

Predictor-development data are collected with participant-level separation.
During controlled intent-elicitation trials, the participant receives a
private spatial task cue that is independent of target ground truth and is not
provided to the autonomous planner. The cue induces a real intended direction
or destination for joystick control without turning that intention into target
evidence. Voice-specified destinations are excluded from predictor training
because they already provide explicit responsibility.

No participant contributes to more than one of training, validation, or final
predictor test. The temporal configuration $(W,H,F_s)$ and feature family
$F_k$ are selected only on training and validation participants. Candidate
horizons are rejected when $H_{\rm eff}\le0$ or when upper-tail FDE makes
responsibility redistribution unstable. The selected temporal configuration,
feature family, RF parameters, and prediction-trigger threshold are frozen
before the held-out predictor test and all mission-level comparisons.

Primary predictor characterization reports median, 90th-percentile, and
95th-percentile FDE. Predefined feature-family ablation tests whether gaze,
view context, and switch context add predictive value beyond control/dynamics.
This is a secondary research question: the main HRI claim does not depend on
gaze improving prediction.

\subsection{Stage E: Collaboration Conditions}

Five conditions isolate adaptation to human intervention and the information
available about human spatial responsibility.

\paragraph{B0: Autonomous}
All USVs remain autonomous. B0 provides an autonomous reference and contains no
human-controlled responsibility center.

\paragraph{B1: Human No-Adapt}
One USV is directly human controlled, but the intervention itself does not
trigger immediate fleet redistribution. Ordinary sensing-, arrival-, and
path-validity replanning remains active.

\paragraph{B2: Human Reactive-Current}
One USV is directly human controlled. The planner represents current human
spatial responsibility by fixing the current human-controlled position,

\begin{equation}
c_h^{B2}(t)=p_h(t).
\label{eq:b2_stagee}
\end{equation}

The remaining autonomous fleet is redistributed using the same constrained
operator as B3.

\paragraph{B3: Human Predictive-Future}
The frozen predictor supplies

\begin{equation}
c_h^{B3}(t)=\hat p_h(t+H)
\label{eq:b3_stagee}
\end{equation}

as the temporary fixed center while direct human control is active.

\paragraph{B4: Oracle Future Responsibility}
For controlled replay analysis only, the actual future human-controlled
position at the same horizon is supplied,

\begin{equation}
c_h^{B4}(t)=p_h(t+H).
\label{eq:b4_stagee}
\end{equation}

B4 never uses target ground truth and is not a live participant condition. It
estimates the upper bound of the same redistribution mechanism under perfect
future-human-position information.

The primary system comparison is B2 versus B3. B1 measures the consequence of
human intervention without proactive fleet adaptation; B4 separates predictor
error from limitations of the redistribution mechanism. Explicit voice
commands are validated separately as an explicit-responsibility capability and
do not redefine B2.

\subsection{Controlled and Live Evaluation}

The evaluation has three complementary layers. First, a controlled Monte Carlo
study samples the frozen calibrated observation model and compares B0--B4 over
matched maps, target positions, initial fleets, and frozen human-control traces.
This isolates responsibility-allocation effects from incidental image-level
perception variability. Second, perception-in-the-loop trials process Gazebo
images through the frozen detector and world-localization pipeline online,
validating the complete image-to-belief-to-planning chain. Third, the live HRI
study compares B1, B2, and B3 with counterbalanced within-subject condition
order. B0 remains an autonomous reference and B4 remains an offline/replay
oracle.

Replay of the same timestamped human trace may be used for paired algorithmic
comparisons, but replay is treated as an experimental-control device rather
than an independent human observation.

\subsection{World Assignment to Experimental Questions}

The calibration world is reserved for sensing identification. Open-search
missions evaluate the basic probabilistic-search loop without strong path
constraints. Channel missions expose path-cost assignment effects. Archipelago
missions provide multiple semi-open spatial responsibility regions and form the
principal environment for predictive human redistribution. Map selection is
fixed before B2--B3 testing and is not tuned to maximize predictive benefit.

\subsection{Primary System Outcome}

The unique primary mission-level outcome is time to calibrated target
confirmation under a fixed mission budget $\tau=T_{\rm budget}$. For condition
$c$, let

\begin{equation}
\mathcal{S}_c(t)=P(T_{\rm conf}>t)
\label{eq:mission_survival}
\end{equation}

be the survival function of confirmation time. Restricted mean
time-to-confirmed-detection is

\begin{equation}
\mathrm{RMST}_c(\tau)=\int_0^{\tau}\mathcal{S}_c(t)\,dt.
\label{eq:rmst}
\end{equation}

The predeclared primary estimand is

\begin{equation}
\Delta_{\rm RMST}
=\mathrm{RMST}_{B3}(\tau)-\mathrm{RMST}_{B2}(\tau),
\label{eq:delta_rmst}
\end{equation}

where $\Delta_{\rm RMST}<0$ favors predictive-future redistribution. This
estimand incorporates right-censored missions instead of averaging only
successful runs.

\subsection{Secondary System and Mechanism Outcomes}

Success probability before budget is secondary rather than co-primary.
Additional secondary system outcomes are total fleet travel distance, maximum
single-USV travel distance, accepted replanning count, planning latency, and
minimum inter-USV separation.

The principal mechanism metric is \emph{Redundant Sensing Overlap} (RSO). For
cell $i$, define the accumulated effective sensing hazard generated by the
human-controlled vehicle and autonomous fleet as

\begin{equation}
H_{h,i}=\int_0^{T_{\rm end}} h_{hi}(t)\,dt,
\label{eq:human_hazard_effort}
\end{equation}

\begin{equation}
H_{A,i}=\int_0^{T_{\rm end}}
\sum_{r\in\mathcal{R}_A(t)}h_{ri}(t)\,dt.
\label{eq:auto_hazard_effort}
\end{equation}

For human-intervention conditions with nonzero human sensing effort,

\begin{equation}
\mathrm{RSO}=
\frac{\sum_i b_i^0\min(H_{h,i},H_{A,i})}
{\sum_i b_i^0 H_{h,i}}.
\label{eq:rso}
\end{equation}

RSO measures the fraction of belief-weighted human sensing effort that is also
covered by the autonomous fleet. It is undefined for B0 and is not replaced by
a geometric trajectory-overlap metric.

\subsection{Human-Centered Outcomes}

The live study includes two secondary HRI outcomes. Perceived workload is
measured using the six standard NASA-TLX dimensions and summarized by the
unweighted Raw TLX score. Coordination alignment is measured with the fixed
1--7 Likert item: ``The autonomous fleet adapted in a way that was consistent
with where I intended to search.'' This single task-specific item is treated as
a diagnostic outcome rather than as a validated psychometric scale.

\subsection{Human-Subject Design}

The final live participant study uses a within-subject design for B1, B2, and
B3. The primary contrast is B2 versus B3; condition order is randomized or
counterbalanced. The same world families, mission budgets, detector, calibrated
sensor model, controller, planner, and frozen predictor are used across
conditions. Target placements and initial fleet configurations are drawn from
a predeclared mission set disjoint from detector development and sensor
calibration.

Participant identity is treated as a repeated factor rather than treating
sliding predictor windows as independent samples. Sample size is determined
before formal data collection from a power or precision analysis based on
pilot participant-level variance for the primary $\Delta_{\rm RMST}$ estimand.
No participant is added or removed based on interim significance.

\subsection{Frozen Hypotheses}

The primary hypothesis is that B3 reduces RMST relative to B2,
$\Delta_{\rm RMST}<0$. The principal mechanism hypothesis is that B3 reduces
RSO relative to B2. B4 is expected to define the attainable upper-bound trend
of the same responsibility mechanism and is used to quantify the gap due to
prediction error. A separate secondary predictor hypothesis tests whether
valid gaze/view context improves held-out future-position prediction beyond
control/dynamics alone. Failure of this secondary predictor hypothesis does not
alter the primary B2--B3 HRI question.

\subsection{Statistical Analysis}

Kaplan--Meier curves are used to visualize time to calibrated confirmation, but
the predeclared primary summary is $\Delta_{\rm RMST}$ at the common budget
$\tau$. The B2--B3 primary effect is reported with an effect estimate and 95\%
confidence interval using a paired participant-level or mixed-effects procedure
consistent with the counterbalancing structure. Success probability is analyzed
as a secondary paired binary or mixed-effects outcome. RSO, travel distance,
replanning, NASA-TLX, and coordination-alignment scores use paired
participant-level contrasts or hierarchical models, with robust intervals when
distributional assumptions are weak.

Predictor analyses remain separate from mission analyses. Participant-level
FDE distributions, feature-family ablation, and prediction-trigger statistics
are reported on held-out participants. Formal mission data are never used to
retrain the predictor, retune the detector, change the B2/B3 definitions,
select the primary metric, or recalibrate the sensor model.

% ------------------------------------------------------------
% FIG. 6: EXPERIMENTAL LOGIC
% Suggested final file: figures/fig6_experimental_design.pdf
% Recommended layout: upper row = freeze/data-isolation pipeline;
% lower row = B0--B4 information/timing comparison.
% ------------------------------------------------------------
\begin{figure*}[t]
\centering
\figplaceholderwide{50mm}{
\textbf{Frozen experimental workflow and HRI condition ladder.}\\
Upper row: detector development/freeze $\rightarrow$ sensor calibration/freeze
$\rightarrow$ Tobii task-space validation $\rightarrow$ participant-disjoint
predictor development/freeze $\rightarrow$ controlled Monte Carlo
$\rightarrow$ perception-in-the-loop $\rightarrow$ live HRI. Mark all data
roles as disjoint.\\
Lower row: B0 autonomous reference; B1 human no-adapt; B2
Reactive-Current with fixed center $p_h(t)$; B3 Predictive-Future with fixed
center $\hat p_h(t+H)$; B4 offline Oracle Future with $p_h(t+H)$. Highlight B2
vs B3 as the primary comparison. Mark live conditions B1/B2/B3 and offline B4.
Add compact metric labels: Primary $\Delta_{\rm RMST}$; mechanism RSO; human
outcomes Raw NASA-TLX and coordination alignment.}
\caption{Frozen experimental workflow and present--predicted--oracle HRI
evaluation ladder. B2 and B3 use the same redistribution operator and differ
only in whether current or predicted future human spatial responsibility is
available.}
\label{fig:experiment_workflow}
\end{figure*}

% ============================================================
% X. RESULTS TEMPLATE
% ============================================================
\section{Results}

The final manuscript inserts only results produced after all model and protocol
freezes described above. This master draft deliberately separates the result
structure from the method so that platform development cannot silently change
a hypothesis after observing mission outcomes.

\subsection{Sensor and Interface Validation}

% FUTURE RESULT FIGURES -- insert only after frozen calibration:
% figures/result_pd_surface.pdf       : calibrated P_D(range,bearing)
% figures/result_localization_error.pdf : localization residuals/anchor choice
% Optional table: P_FA, nu_eff, and confirmation false-termination statistics.

The final paper reports the calibrated $P_D(\rho,\beta)$ model and uncertainty,
$P_{\rm FA}$, localization residual statistics and selected anchor, temporal
model and $\nu_{\rm eff}$, the frozen confirmation rule, Tobii task-space error,
valid-sample rate, and cross-host timing residual. These values establish the
measurement validity of the subsequent human study.

\subsection{Predictor Development and Held-Out Performance}

% FUTURE RESULT FIGURES -- insert only after participant-disjoint testing:
% figures/result_fde_horizon.pdf      : median/Q90/Q95 FDE versus horizon
% figures/result_feature_ablation.pdf : predeclared F0--F4 ablation

The final paper reports the selected $(W,H,F_s)$, selected feature family,
held-out participant FDE distribution, effective anticipation time, and
predefined feature-family ablation. The purpose is to determine whether gaze/view context adds predictive value
beyond control/dynamics and whether the selected horizon is useful to the fleet
rather than merely easy to predict.

\subsection{Mission-Level Collaboration Results}

% FUTURE RESULT FIGURES -- insert only after formal frozen experiments:
% figures/result_survival_b0_b4.pdf : time-to-confirmation survival curves
% figures/result_b2_b3_rmst.pdf     : paired/participant-level Delta_RMST
% figures/result_rso.pdf            : redundant sensing overlap
% figures/result_hri_outcomes.pdf   : NASA-TLX and coordination alignment

The primary result is the B2--B3 difference in RMST. Success probability,
RSO, fleet effort, replanning, and safety are secondary system or mechanism
outcomes. The live HRI analysis additionally reports Raw NASA-TLX and the
predeclared coordination-alignment item for B1--B3. B0 contextualizes the
autonomous reference, while B4 is reported only in controlled/replay analyses
as an oracle future-responsibility bound. Controlled Monte Carlo,
perception-in-the-loop, and live participant results are reported separately
rather than pooled.

\begin{table*}[t]
\centering
\caption{Frozen mission-result table template. Values are inserted only after
the predeclared analyses are completed.}
\label{tab:result_template}
\small
\setlength{\tabcolsep}{4pt}
\begin{tabular}{lcccccc}
\toprule
Condition & RMST (primary) & Success & RSO & Fleet distance & Replans & Min. sep.\\
\midrule
B0 Autonomous & \TBD{result} & \TBD{result} & N/A & \TBD{result} & \TBD{result} & \TBD{result}\\
B1 Human No-Adapt & \TBD{result} & \TBD{result} & \TBD{result} & \TBD{result} & \TBD{result} & \TBD{result}\\
B2 Reactive-Current & \TBD{result} & \TBD{result} & \TBD{result} & \TBD{result} & \TBD{result} & \TBD{result}\\
B3 Predictive-Future & \TBD{result} & \TBD{result} & \TBD{result} & \TBD{result} & \TBD{result} & \TBD{result}\\
B4 Oracle Future$^{\dagger}$ & \TBD{result} & \TBD{result} & \TBD{result} & \TBD{result} & \TBD{result} & \TBD{result}\\
\bottomrule
\end{tabular}
\\
\footnotesize{$^{\dagger}$B4 is an offline/replay oracle, not a live participant condition.}
\end{table*}

\begin{table}[t]
\centering
\caption{Human-centered result template for live B1--B3 conditions.}
\label{tab:hri_result_template}
\footnotesize
\setlength{\tabcolsep}{3pt}
\begin{tabular}{lcc}
\toprule
Condition & Raw TLX & Alignment\\
\midrule
B1 No-Adapt & \TBD{result} & \TBD{result}\\
B2 Reactive-Current & \TBD{result} & \TBD{result}\\
B3 Predictive-Future & \TBD{result} & \TBD{result}\\
\bottomrule
\end{tabular}
\end{table}

% ============================================================
% XI. DISCUSSION
% ============================================================
\section{Discussion}

The central HRI contribution is the distinction between uncertainty about the
target and uncertainty about the human's future contribution. The Bayesian
posterior represents target uncertainty; human spatial responsibility
represents where human search effort is expected to occur. This separation
prevents joystick behavior, gaze, voice, or predicted motion from being
misinterpreted as target evidence while still allowing the autonomous fleet to
adapt to human intervention.

The primary comparison is intentionally between B2 and B3. B2 already adapts
to the human by fixing the current human-controlled position $p_h(t)$, so B3
cannot gain simply because one condition has a human center and the other does
not. Any additional B3 benefit must arise from actionable information about
future spatial responsibility. B4 then quantifies the remaining gap under
perfect future-human-position information without introducing target truth.

The fixed-center formulation does not manually delete a region that the human
is expected to search. All searchable free cells remain in the same
responsibility optimization. This makes overlap possible when the belief and
geometry justify it, while RSO measures whether autonomous sensing unnecessarily
repeats human sensing effort.

Prediction is useful only when it provides meaningful coordination lead time.
A highly accurate short-horizon predictor may offer little benefit if planning
and motion consume the available horizon, while a long horizon may provide more
lead time at the cost of larger endpoint uncertainty. FDE, $H_{\rm eff}$, RSO,
and the B2--B3 RMST contrast therefore evaluate prediction in the HRI role that
matters: whether the fleet can coordinate earlier and more effectively with the
human.

The interaction modalities are deliberately asymmetric. Joystick control
creates implicit future responsibility, gaze may add predictive context, and
voice provides explicit vehicle--destination responsibility. The paper does not
claim a new multimodal fusion algorithm, and the primary HRI hypothesis remains
valid even if gaze does not improve the held-out predictor. Likewise, the
implemented voice channel should not be described as general high-level intent
understanding.

A natural extension is to represent richer human task semantics. Future work
could map natural-language instructions such as search a region, hold position,
patrol, pursue, escort, or form up into structured fleet responsibilities,
potentially using large language or multimodal models. That extension would
change the semantic interpreter, but the present responsibility interface could
remain the downstream contract. The current study intentionally restricts
voice to explicit vehicle--destination responsibility in order to isolate the
spatial-coordination mechanism.

Static LOS qualification is conservative: it prevents known opaque static
obstacles from producing unjustified negative evidence, but it is not a full
3-D occlusion model. Dynamic vessel occlusion and transient wave effects remain
part of residual sensing uncertainty. The current allocation objective also
uses Euclidean belief-weighted geometry for center generation and obstacle-aware
path length only during assignment. This interpretable separation is retained
to keep the HRI hypothesis identifiable; path-aware center generation is a
possible extension if obstacle topology is later shown to dominate the human
responsibility effect.

% ============================================================
% XII. CONCLUSION
% ============================================================
\section{Conclusion}

This paper formulates collaborative multi-USV probabilistic search as an HRI
problem of adapting autonomous search responsibility to human intervention.
Calibrated visual observations update the Bayesian target belief, while human
interaction is represented separately as spatial responsibility. Under direct
human control, B2 uses the current human-controlled position $p_h(t)$ as the
reactive responsibility center, whereas B3 uses the predicted future position
$\hat p_h(t+H)$. An offline B4 oracle uses $p_h(t+H)$ to quantify the
performance gap attributable to prediction error. All three conditions share
the same fixed-center redistribution operator.

The proposed mechanism keeps all searchable free cells active, recomputes only
autonomous centers around fixed human responsibility, and then executes the
allocation through pairwise obstacle-aware A* planning and path-cost Hungarian
assignment. Voice provides a separate explicit vehicle--destination
responsibility channel, while gaze is treated as optional predictive context;
neither is allowed to modify the target posterior without an explicit human
observation model.

The experimental contract freezes the primary HRI question before mission
results are observed. The unique primary system estimand is the B2--B3
difference in restricted mean time-to-confirmed-detection. RSO tests the
redundant-sensing mechanism, while success probability, fleet effort, Raw
NASA-TLX, and perceived coordination alignment provide secondary system and
human-centered evidence. The study therefore asks a focused question: whether
anticipating the future spatial consequence of human action enables a robot
fleet to cooperate with the human earlier and more effectively than reacting to
where the human-controlled vehicle is now.

% ============================================================
% REFERENCES
% ============================================================
\begin{thebibliography}{00}

\bibitem{b1}
F. Bourgault, A. Goktogan, T. Furukawa, and H. F. Durrant-Whyte,
``Coordinated search for a lost target in a Bayesian world,''
\emph{Advanced Robotics}, vol. 18, no. 10, pp. 979--1000, 2004.

\bibitem{b2}
D. Xie, Y. Lin, R. A. Grupen, and A. R. Hanson,
``Intention-based coordination and interface design for human-robot cooperative
search,'' in \emph{Proc. IEEE/RSJ Int. Conf. Intelligent Robots and Systems
(IROS)}, pp. 3349--3354, 2011.

\bibitem{b3}
L. Heintzman, A. Hashimoto, N. Abaid, and R. K. Williams,
``Anticipatory planning and dynamic lost person models for human-robot search
and rescue,'' in \emph{Proc. IEEE Int. Conf. Robotics and Automation (ICRA)},
pp. 8252--8258, 2021.

\bibitem{b4}
G. Hoffman, T. Bhattacharjee, and S. Nikolaidis,
``Inferring human intent and predicting human action in human--robot
collaboration,'' \emph{Annual Review of Control, Robotics, and Autonomous
Systems}, vol. 7, pp. 73--95, 2024.

\bibitem{b5}
M. Dalmasso, J. E. Dom\'inguez-Vidal, I. J. Torres-Rodr\'iguez,
P. Jim\'enez, A. Garrell, and A. Sanfeliu,
``Shared task representation for human--robot collaborative navigation: The
collaborative search case,'' \emph{International Journal of Social Robotics},
vol. 16, pp. 145--171, 2024.

\bibitem{b7}
P. M. Dames,
``Distributed multi-target search and tracking using the PHD filter,''
\emph{Autonomous Robots}, vol. 44, pp. 673--689, 2020.

\bibitem{b8}
F. Jovan, M. Tomy, N. Hawes, and J. Wyatt,
``Efficiently exploring for human robot interaction: partially observable
Poisson processes,'' \emph{Autonomous Robots}, vol. 47, pp. 121--138, 2023.

\end{thebibliography}

\end{document}